\section{Discussion}
\label{sec:discussion}

\subsection{Summary of Problem and Results}

This paper modeled series RLC circuits using first- and second-order ODEs,
deriving analytic solutions for RC, RL, and RLC configurations.
The resistor was identified as the sole energy-dissipating element, and its
value shown to directly determine the discriminant of the characteristic
equation~\eqref{eq:char}, thereby controlling whether the circuit is
overdamped, critically damped, or underdamped.

\subsection{Interpretation in Terms of the Research Question}

Without resistance ($R = 0$), the characteristic roots are purely imaginary
($r = \pm j\omega_0$), yielding perpetual, undamped oscillations in an ideal
LC tank circuit.
Any nonzero $R$ introduces the decay envelope $e^{-\alpha t}$ that eventually
damps oscillations to zero.
In engineering applications: bandpass filters rely on controlled underdamping
to achieve frequency selectivity when the output is taken across the resistor,
while the capacitor-voltage response studied here behaves like a low-pass
response with possible resonant overshoot. Overdamped circuits ensure no
ringing in power supplies and ADC front-ends; critically damped designs appear
in servo systems where settling speed is paramount. A passive RLC circuit can
store and exchange energy, but a sustained oscillator also requires an active
energy source or feedback element to replace resistive losses.
This analytical framework directly addresses the research question: resistors
are essential for stabilizing oscillatory behavior and enabling practical
circuit design.

\subsection{Limitations and Future Work}

The lumped-element and linearity assumptions restrict the model's validity to
low frequencies and small signals.
Future work should incorporate:
(i) a nonlinear resistance model $R(T,i)$ to capture self-heating;
(ii) distributed-element analysis for RF applications above \qty{1}{\mega\hertz};
(iii) experimental validation using an oscilloscope to compare measured
$\vC(t)$ waveforms against the analytic solutions derived here; and
(iv) extension to parallel RLC and more complex multi-mesh circuits modeled
by systems of ODEs.
